\documentclass[11pt]{article}
\usepackage[margin=2.5cm]{geometry}
\usepackage{amsmath,amssymb,bm}
\usepackage{booktabs}
\usepackage{hyperref}

\title{A Generalized Framework for Direct Visual Servoing\\with Gaussian-based Scene Representations}
\author{Youssef Alj \and Guillaume Caron}
\date{Working Draft --- April 2026}

\begin{document}
\maketitle

\section{Generalized Framework}

\subsection{Photometric Visual Servoing: Unified Formulation}

The visual servoing task minimizes the Sum of Squared Differences (SSD) between visual features extracted from the current and desired images:
\begin{equation}
    \mathcal{C}(\mathbf{p}) = \frac{1}{2} \| \mathbf{s}(\mathbf{p}) - \mathbf{s}^* \|^2
\end{equation}
where $\mathbf{p} \in \mathbb{R}^n$ is the camera pose ($n=6$ for eye-in-hand), $\mathbf{s}$ is the visual feature vector, and $\mathbf{s}^*$ is the desired feature.

The control law is based on the Gauss-Newton (or Levenberg-Marquardt) optimization:
\begin{equation}
    \mathbf{v} = -\lambda \left( \mu \, \text{diag}(\mathbf{H}) + \mathbf{H} \right)^{-1} \mathbf{L}_s^\top (\mathbf{s} - \mathbf{s}^*)
\end{equation}
where $\mathbf{L}_s$ is the interaction matrix, $\mathbf{H} = \mathbf{L}_s^\top \mathbf{L}_s$ is the approximate Hessian, $\lambda$ is the gain, and $\mu$ is the LM damping factor.

\subsection{The Convergence Domain Problem}

PVS achieves sub-millimeter precision at convergence but suffers from a narrow convergence domain due to the highly nonlinear cost function. Prior work has shown that \emph{smoothing} the visual features enlarges the convergence basin at the expense of precision:

\begin{center}
\renewcommand{\arraystretch}{1.3}
\begin{tabular}{lccl}
\toprule
\textbf{Method} & \textbf{Smoothing Domain} & \textbf{Feature} & \textbf{Key Parameter} \\
\midrule
PVS~\cite{collewet2011} & None & $\mathbf{I}(\mathbf{u})$ & --- \\
PGM-VS~\cite{crombez2019tro} & 2D image space & $G(\mathbf{I}, \lambda_g)$ & $\lambda_g$ (Gaussian spread) \\
DDVS~\cite{caron2021ral} & Optical (thin lens) & $I_d(\mathbf{u})$ & $D, Z_f$ (aperture, focus) \\
\textbf{Ours (3DGS-VS)} & \textbf{3D scene space} & $\mathbf{I}_\alpha(\mathbf{u})$ & $\alpha$ (scale factor) \\
\bottomrule
\end{tabular}
\end{center}

\subsection{Unifying Through Gaussian Spread}

All three smoothing-based methods can be understood through a common lens: they control the \emph{effective Gaussian spread} $\sigma_{\text{eff}}$ in the image plane, which determines the convergence domain radius.

\subsubsection{PGM-VS: 2D Isotropic Gaussian Convolution}

PGM-VS~\cite{crombez2019tro} replaces pixel intensities with Photometric Gaussian Mixtures:
\begin{equation}
    G(\mathbf{I}, \mathbf{u}_g, \lambda_g) = \sum_{\mathbf{u}} I(\mathbf{u}) \exp\left( -\frac{\|\mathbf{u}_g - \mathbf{u}\|^2}{2\lambda_g^2} \right)
\end{equation}
The spread $\sigma_{\text{eff}} = \lambda_g$ is uniform across the image (isotropic, depth-independent). The interaction matrix $\mathbf{L}_G$ is derived analytically from the Gaussian derivatives $K_u$, $K_v$.

\textbf{Properties:}
\begin{itemize}
    \item Large $\lambda_g$ $\Rightarrow$ wide convergence domain, low precision
    \item Small $\lambda_g$ $\Rightarrow$ narrow domain, high precision
    \item Multi-level strategy: $\lambda_g^{(1)} > \lambda_g^{(2)} > \ldots > \lambda_g^{(L)}$
    \item Complexity: $O(NM)$ with separable convolution (TRO 2019 model)
\end{itemize}

\subsubsection{DDVS: Optical Defocus Blur}

DDVS~\cite{caron2021ral} exploits the thin lens camera model. A 3D point $\mathbf{X}$ at depth $Z$ projects to a Circle of Confusion (CoC) with diameter:
\begin{equation}
    d(Z) = \frac{Df}{Z_f - f} \left( 1 - \frac{Z_f}{Z} \right)
\end{equation}
The effective Gaussian spread is $\sigma_{\text{eff}}(Z) = d(Z) / (6 k_u)$, which is \emph{depth-dependent} and controlled by the aperture $D$ and focus depth $Z_f$.

\textbf{Properties:}
\begin{itemize}
    \item Blur is anisotropic and depth-dependent
    \item Directly acquired (no image processing needed)
    \item Complexity: $O(NM)$ (same as PVS)
    \item Requires physical control of camera aperture/focus
\end{itemize}

\subsubsection{3DGS-VS (Ours): 3D Gaussian Scale Inflation}

In 3D Gaussian Splatting, the scene is represented as a mixture of 3D Gaussians:
\begin{equation}
    \mathbf{I}(\mathbf{u}) = \sum_{k=1}^{K} c_k \, \alpha_k \prod_{j<k}(1-\alpha_j) \, \mathcal{N}(\mathbf{u}; \boldsymbol{\mu}_k^{2D}, \boldsymbol{\Sigma}_k^{2D})
\end{equation}
where $\boldsymbol{\mu}_k^{2D}$ and $\boldsymbol{\Sigma}_k^{2D}$ are the projected 2D mean and covariance of the $k$-th Gaussian.

We scale the 3D Gaussian covariances by a factor $\alpha$:
\begin{equation}
    \boldsymbol{\Sigma}_k^{3D}(\alpha) = \alpha^2 \boldsymbol{\Sigma}_k^{3D}
\end{equation}
The projected 2D covariance becomes:
\begin{equation}
    \boldsymbol{\Sigma}_k^{2D}(\alpha) = \mathbf{J}_k \, \alpha^2 \boldsymbol{\Sigma}_k^{3D} \, \mathbf{J}_k^\top
\end{equation}
where $\mathbf{J}_k$ is the Jacobian of the projection at the $k$-th Gaussian's mean.

The effective spread in the image plane is:
\begin{equation}
    \sigma_{\text{eff},k}(\alpha) = \alpha \cdot \sigma_k^{2D}
\end{equation}
which is \emph{per-Gaussian, depth-dependent, and anisotropic} --- inheriting the 3D structure of the scene.

\textbf{Properties:}
\begin{itemize}
    \item Blur is native to the 3D scene representation (not a post-processing step)
    \item Depth-dependent and anisotropic by construction
    \item Each Gaussian's contribution is scaled independently based on its 3D geometry
    \item Complexity: $O(NM)$ (just re-rendering with modified scales, same pipeline)
    \item No physical camera modification needed (virtual rendering)
    \item Compatible with any camera model supported by gsplat (pinhole, fisheye, etc.)
\end{itemize}

\subsubsection{Impact of $\alpha$ in the 3DGS Rendering Equation}

We now trace the effect of $\alpha$ through the complete 3DGS rendering pipeline to show its mathematical impact on the rendered image and, consequently, on the cost function.

\paragraph{3D Gaussian representation.}
Each Gaussian $k$ is parameterized by its mean $\boldsymbol{\mu}_k \in \mathbb{R}^3$, rotation $\mathbf{R}_k \in SO(3)$, scales $\mathbf{s}_k = (s_x, s_y, s_z)_k$, opacity $o_k$, and spherical harmonics coefficients $\mathbf{c}_k$. The 3D covariance matrix is:
\begin{equation}
    \boldsymbol{\Sigma}_k = \mathbf{R}_k \, \mathbf{S}_k \mathbf{S}_k^\top \, \mathbf{R}_k^\top, \quad \mathbf{S}_k = \text{diag}(s_x, s_y, s_z)_k
\end{equation}
Inflating by $\alpha$: $\mathbf{S}'_k = \alpha \, \mathbf{S}_k$, yielding:
\begin{equation}
    \boldsymbol{\Sigma}'_k = \alpha^2 \, \boldsymbol{\Sigma}_k
    \label{eq:cov3d_inflated}
\end{equation}

\paragraph{Splatting (3D to 2D projection).}
The EWA (Elliptical Weighted Average) splatting projects each 3D Gaussian onto the image plane. Given a camera with view matrix $\mathbf{W}$ and the Jacobian $\mathbf{J}_k$ of the local affine approximation of the projection at $\boldsymbol{\mu}_k$:
\begin{equation}
    \boldsymbol{\Sigma}_{k}^{2D} = \mathbf{J}_k \, \mathbf{W} \, \boldsymbol{\Sigma}_k \, \mathbf{W}^\top \mathbf{J}_k^\top
\end{equation}
With the inflated covariance~\eqref{eq:cov3d_inflated}:
\begin{equation}
    \boldsymbol{\Sigma}_{k}^{\prime 2D} = \alpha^2 \, \boldsymbol{\Sigma}_{k}^{2D}
    \label{eq:cov2d_inflated}
\end{equation}
Note that $\mathbf{J}_k$ and $\mathbf{W}$ are \emph{unchanged} --- the inflation affects only the covariance, not the Gaussian centers or the camera pose.

\paragraph{Per-pixel alpha evaluation.}
The contribution of Gaussian $k$ at pixel $\mathbf{p}$ is evaluated via the 2D Gaussian:
\begin{equation}
    \alpha_k(\mathbf{p}) = o_k \cdot \exp\!\left( -\frac{1}{2} \, \mathbf{d}_k^\top \left(\boldsymbol{\Sigma}_{k}^{2D}\right)^{-1} \mathbf{d}_k \right)
\end{equation}
where $\mathbf{d}_k = \mathbf{p} - \boldsymbol{\mu}_k^{2D}$ is the offset from the projected mean. Substituting the inflated covariance~\eqref{eq:cov2d_inflated}:
\begin{equation}
    \alpha'_k(\mathbf{p}) = o_k \cdot \exp\!\left( -\frac{1}{2\alpha^2} \, \mathbf{d}_k^\top \left(\boldsymbol{\Sigma}_{k}^{2D}\right)^{-1} \mathbf{d}_k \right)
    \label{eq:alpha_inflated}
\end{equation}
\textbf{The exponent is divided by $\alpha^2$.} This has two direct consequences:
\begin{enumerate}
    \item The Gaussian's footprint in the image plane grows by factor $\alpha$ (each pixel sees contributions from Gaussians that are $\alpha$ times farther away).
    \item For a pixel at distance $\mathbf{d}_k$ from the Gaussian center, $\alpha'_k > \alpha_k$, meaning distant Gaussians contribute \emph{more} to each pixel --- a smoothing effect.
\end{enumerate}

\paragraph{Alpha compositing (final pixel color).}
The rendered color at pixel $\mathbf{p}$ follows front-to-back alpha compositing:
\begin{equation}
    C(\mathbf{p}) = \sum_{k=1}^{K} c_k \, \alpha'_k(\mathbf{p}) \prod_{j < k} \left(1 - \alpha'_j(\mathbf{p})\right)
    \label{eq:compositing}
\end{equation}
Since $\alpha'_k(\mathbf{p})$ is affected by $\alpha$ via~\eqref{eq:alpha_inflated}, both the direct contribution and the occlusion term $\prod(1-\alpha'_j)$ are modified. Specifically:
\begin{itemize}
    \item Foreground Gaussians become more transparent (their $\alpha'_k$ decreases at their center when $\alpha > 1$ because the same mass is spread over a wider area).
    \item Background Gaussians become more visible through the foreground.
    \item The overall image is a \emph{weighted average} of more Gaussians per pixel, producing a blurred appearance.
\end{itemize}

\paragraph{Effect on image gradients.}
Let $\nabla_{\mathbf{u}} C$ denote the spatial gradient of the rendered image. Since the compositing~\eqref{eq:compositing} involves smoother per-Gaussian evaluations (wider support), the resulting image gradient magnitude decreases:
\begin{equation}
    \| \nabla_{\mathbf{u}} C_\alpha \| \lesssim \frac{1}{\alpha} \| \nabla_{\mathbf{u}} C_1 \|
\end{equation}
This is analogous to the classical result in scale-space theory: convolving with a Gaussian of width $\sigma$ reduces gradient magnitude by $O(1/\sigma)$.

\paragraph{Effect on the photometric cost function.}
The photometric cost at a perturbed pose $\Delta\mathbf{p}$ from the desired pose is:
\begin{equation}
    \mathcal{C}_\alpha(\Delta\mathbf{p}) = \frac{1}{2} \sum_{\mathbf{u}} \left( C_\alpha(\mathbf{u}; \Delta\mathbf{p}) - C_\alpha^*(\mathbf{u}) \right)^2
\end{equation}
Since $C_\alpha$ is smoother, $\mathcal{C}_\alpha$ inherits the smoothness:
\begin{itemize}
    \item Local minima are suppressed (small-scale cost variations are averaged out).
    \item The main basin around $\Delta\mathbf{p} = 0$ widens.
    \item The gradient $\nabla_{\Delta\mathbf{p}} \mathcal{C}_\alpha$ is weaker but more directionally consistent over a larger region.
\end{itemize}
This explains the experimental observation of an \textbf{optimal $\alpha$}: too small, and local minima persist; too large, and the gradient vanishes entirely, preventing convergence.

\paragraph{Comparison with PGM-VS.}
In PGM-VS, the smoothing is applied \emph{after} rendering via a 2D Gaussian convolution:
\begin{equation}
    G(\mathbf{I}, \mathbf{u}_g, \lambda_g) = \sum_{\mathbf{u}} I(\mathbf{u}) \exp\!\left(-\frac{\|\mathbf{u}_g - \mathbf{u}\|^2}{2\lambda_g^2}\right)
\end{equation}
This is equivalent to convolving the rendered image with an isotropic Gaussian kernel of width $\lambda_g$. The crucial difference with our approach:

\begin{center}
\renewcommand{\arraystretch}{1.3}
\begin{tabular}{lcc}
\toprule
 & \textbf{PGM-VS} & \textbf{3DGS-VS (ours)} \\
\midrule
Smoothing stage & After rendering (2D) & Before rendering (3D) \\
Depth-dependence & No ($\lambda_g$ uniform) & Yes ($\sigma_k^{2D} \propto 1/Z_k$) \\
Anisotropy & No (isotropic in 2D) & Yes (inherits 3D shape) \\
Occlusion handling & Ignores (post-rendering blur) & Respects (compositing order) \\
Interaction matrix & Requires $K_u, K_v$ convolutions & Standard PVS $\mathbf{L}_I$ \\
\bottomrule
\end{tabular}
\end{center}

Our method's key advantage is that the smoothing \emph{respects the 3D scene geometry}: a foreground object and a background wall are blurred independently according to their depth and 3D shape, rather than being smeared together by a uniform 2D convolution. This produces a more physically meaningful smoothed image and avoids spurious cost function minima caused by blending unrelated scene regions.

\subsection{Connection to Convergence Domain Theory}

Naamani et al.~\cite{naamani2024iros} proved that for DVS with Gaussian blur, the convergence domain radius $r$ satisfies:
\begin{equation}
    r \propto \sigma_{\text{eff}}
\end{equation}
for planar translation (Eq.~18 in~\cite{naamani2024iros}). Our experimental results (Section~\ref{sec:results}) confirm this relationship for 3DGS-VS:
\begin{itemize}
    \item $\alpha = 1.0$ (original): baseline convergence domain
    \item $\alpha = 1.5$: convergence domain approximately doubles for $t_x$
    \item $\alpha = 2.0$: rotation domains increase by $2\times$--$3\times$
    \item $\alpha = 3.0$: gradients vanish, convergence collapses
\end{itemize}

The non-monotonic behavior (convergence improves then collapses) distinguishes 3DGS-VS from the 2D theoretical prediction, which assumes monotonic improvement. This is because 3D scaling introduces depth-dependent effects not captured by the 2D model: different parts of the scene blur at different rates depending on their depth, potentially creating new local minima at extreme scale factors.

\subsection{Scale Adaptation Strategies}

Given that large $\alpha$ helps approach but hurts precision, we consider three adaptation strategies:

\subsubsection{Constant Inflation}

Use a fixed $\alpha > 1$ throughout servoing:
\begin{equation}
    \alpha(t) = \alpha_0, \quad \forall t
\end{equation}
\textbf{Advantage:} Simple, wide basin. \textbf{Disadvantage:} Reduced precision at convergence.

\subsubsection{Coarse-to-Fine}

Decay $\alpha$ over $L$ phases:
\begin{equation}
    \alpha^{(\ell)} = \alpha_0^{1 - \ell/(L-1)}, \quad \ell = 0, 1, \ldots, L-1
\end{equation}
\textbf{Advantage:} Wide basin for approach, full precision for final convergence. \textbf{Disadvantage:} Phase transitions can cause instability; schedule is fixed regardless of convergence state.

\subsubsection{Error-Adaptive}

Scale proportional to normalized error:
\begin{equation}
    \alpha(t) = 1 + (\alpha_0 - 1) \cdot \min\left(\frac{e(t)}{e(0)}, 1\right)
\end{equation}
\textbf{Advantage:} Automatic, no schedule to tune. \textbf{Disadvantage:} Non-monotonic error can keep $\alpha$ high; sensitive to initial error calibration.

\subsection{Comparison Summary}

\begin{center}
\renewcommand{\arraystretch}{1.3}
\begin{tabular}{lccccc}
\toprule
& \textbf{PGM-VS} & \textbf{DDVS} & \textbf{3DGS-VS} & \textbf{3DGS C2F} \\
\midrule
Smoothing domain & 2D image & Optical & 3D scene & 3D scene \\
Depth-dependent & No & Yes & Yes & Yes \\
Anisotropic & No & Yes & Yes & Yes \\
Requires 3D model & No & No & Yes (GS) & Yes (GS) \\
Requires real camera & Yes & Yes & No & No \\
Complexity & $O(NM)$ & $O(NM)$ & $O(NM)$ & $O(NM)$ \\
Convergence speed & Fast & --- & Moderate & Moderate \\
Convergence basin & Moderate & Moderate & Wide & Wide \\
Final precision & High & High & Low & High \\
\bottomrule
\end{tabular}
\end{center}

\section{Experimental Validation}
\label{sec:results}

Our experiments confirm the theoretical framework on four mip-NeRF 360 scenes (room, kitchen, garden, bicycle):

\begin{enumerate}
    \item \textbf{Cost function landscape:} Scaling Gaussians smooths the cost function and widens the basin (Fig.~2 of experimental notes). $\alpha = 1.5$--$2.0$ is optimal; $\alpha = 3.0$ flattens gradients.

    \item \textbf{Convergence domain:} Bisection measurements show $\alpha = 1.5$ doubles translation domain, $\alpha = 2.0$ triples rotation domain. Relationship is non-monotonic (unlike 2D theory).

    \item \textbf{Multi-scene evaluation:} Inflated 3DGS consistently outperforms original PVS across scenes (+20--40\% success rate on medium perturbations). Same $\alpha = 1.8$ works across scenes.

    \item \textbf{PGM-VS comparison:} PGM is 2--8$\times$ faster when it converges but has narrower basin. Inflated 3DGS succeeds on pairs where PGM fails.

    \item \textbf{Qualitative results:} 35 pairs tested (indices 0--150). Inflated converges on 22/35, PGM on 15/35, original on 19/35. The methods are complementary.
\end{enumerate}

\section{Toward a Unified Blur-Adaptive Visual Servoing}

The framework suggests a natural hierarchy for practical systems:

\begin{enumerate}
    \item \textbf{Far from goal} ($\|\mathbf{e}\| \gg 0$): Use inflated 3DGS ($\alpha \approx 2.0$) for maximum convergence basin. The camera navigates coarsely toward the goal.

    \item \textbf{Approaching goal} ($\|\mathbf{e}\|$ moderate): Reduce $\alpha$ progressively (coarse-to-fine) or switch to PGM-VS (which has a smaller basin but faster convergence).

    \item \textbf{Near goal} ($\|\mathbf{e}\| \approx 0$): Use original PVS ($\alpha = 1.0$) or PGM-VS with small $\lambda_g$ for maximum precision.
\end{enumerate}

This hybrid strategy leverages each method's strength: 3DGS scaling for wide-range navigation, PGM for fast approach, and original PVS for final positioning.


\begin{thebibliography}{9}

\bibitem{collewet2011}
C.~Collewet and E.~Marchand,
``Photometric visual servoing,''
\emph{IEEE Trans. Robot.}, vol.~27, no.~4, pp.~828--834, 2011.

\bibitem{crombez2019tro}
N.~Crombez, E.~M.~Mouaddib, G.~Caron, and F.~Chaumette,
``Visual servoing with photometric Gaussian mixtures as dense feature,''
\emph{IEEE Trans. Robot.}, vol.~35, no.~1, pp.~49--63, 2019.

\bibitem{caron2021ral}
G.~Caron,
``Defocus-based direct visual servoing,''
\emph{IEEE Robot. Autom. Lett.}, vol.~6, no.~2, pp.~4057--4064, 2021.

\bibitem{naamani2024iros}
M.~B.~Naamani, G.~Caron, M.~Morisawa, and E.~M.~Mouaddib,
``A mathematical characterization of the convergence domain for direct visual servoing,''
in \emph{Proc. IEEE/RSJ IROS}, pp.~4846--4853, 2024.

\bibitem{guerbas2021}
S.~E.~Guerbas, N.~Crombez, G.~Caron, and E.~M.~Mouaddib,
``Photometric Gaussian Mixtures for direct virtual visual servoing of omnidirectional camera,''
in \emph{IEEE CVPR 2021 Workshop on 3D Vision and Robotics}, 2021.

\end{thebibliography}

\end{document}
